% !TeX root = ./../main.tex

\section{Szum i błędy w komputerach kwantowych}\label{sec:tb:noise}
\agentinfo{
Stan (PHD-TS-95, 2026-11-27): podrozdział skrócony i poprawiony po recenzji \agentref{eksperymenty/recenzje/rozdzial-2-sekcje-2.1-2.2.md}; lista cięć (w tym rysunek z kulą Blocha, równanie Lindblada, miary QV i EPLG, tomografia zbioru bramek, wzór zaniku w teście RB) w \agentref{eksperymenty/recenzje/rozdzial-2-skracanie.md}.
Tabela kanałów tab:tb:noise:channels jest przeniesiona do dodatku \agentref{chapters/appendix_results_tb.tex} (sekcja app:results:tb); wiersz input tego pliku w appendix_results.tex dopisuje orkiestrator.
Etykiety używane przez inne pliki (nie usuwać): etykiety podsekcji, eq:tb:noise:affine i tab:tb:noise:channels (rozdział 8), eq:tb:noise:shrink (rozdział 7), eq:tb:noise:readout (podrozdział 2.3, dodatek A), eq:tb:noise:fringe (rozdziały 3 i 4), eq:tb:noise:ptm-blocks (plik archiwalny noise_structure.tex).
Notacja: parametry kanałów to q_d, q_a, q_phi (nie gamma, bo gamma to funkcja referencyjna), kąt błędu koherentnego to vartheta (nie epsilon, bo epsilon to funkcja PDU), rozkład na zaszumionym urządzeniu to makro probnoisy, macierz odczytu to makro confmat.
Kanał depolaryzujący ma dwie parametryzacje: lambda (Qiskit Aer) i q_d (PennyLane); dla jednego kubitu lambda = 4 q_d / 3.
Model warstwowy: szum po każdym bloku ansatzu, tak jak w rozdziałach 7 i 8 (słowo warstwa pochodzi z kodu).
Macierz transferu Pauliego, jej postać blokowa i liczba obwodów tomografii procesu służą tylko akapitowi par:lr:noise:generators (sekcja 3.5); nazwę wiersz afiniczny usunięto, bo sekcja 3.5 jej nie używa (ustalenie B1 recenzji dawnego rozdziału 9 dotyczy już tylko sekcji 3.5).
Liczby dla ibm_fez pochodzą z migawki kalibracyjnej z 2025-02-26 dołączonej do pakietu qiskit-ibm-runtime 0.49.0, a nie z dni eksperymentów; wartości z dni eksperymentów należą do dodatku \agentref{chapters/appendix_experiments.tex}.
Twierdzenie o układach fotonicznych sprawdzono tylko na poziomie streszczeń.
}

Rzeczywisty procesor kwantowy nie jest układem izolowanym o~dokładnych bramkach i~bezbłędnym pomiarze, jak założono w~podrozdziale~\ref{sec:tb:quantum}.
W~pracy \emph{szumem} nazywa się procesy fizyczne i~niedoskonałości sterowania zaburzające stan rejestru, a~\emph{błędem} ich skutek, czyli różnicę między operacją lub wynikiem otrzymanym a~idealnym.

\subsection{Źródła błędów}\label{subsec:tb:noise:sources}

Podstawowym źródłem błędów jest \gls{decoherence}, czyli utrata kwantowych własności układu wskutek niekontrolowanego oddziaływania z~otoczeniem~\cite{zurek_decoherence_2003,bib:quantum-decoherence}.
W~kubitach nadprzewodzących opisują ją czas relaksacji $\Tone$, z~którym zanika obsadzenie stanu $\ket{1}$, i~czas koherencji $\Ttwo$, z~którym zanikają koherencje (podrozdział~\ref{subsec:tb:density}), a~wraz z~nimi dane zapisane w~fazie~\cite{krantz_quantum_2019}; fazę niszczy zarówno relaksacja, jak i~czyste defazowanie o~stałej czasowej $\Tphi$, stąd $1/\Ttwo=1/(2\Tone)+1/\Tphi$.
Pozostałe źródła to błędy sterowania~\cite{krantz_quantum_2019}, błędy przygotowania stanu i~pomiaru (SPAM), \gls{crosstalk}, czyli zaburzanie kubitu przez operacje wykonywane w~jego sąsiedztwie~\cite{sarovar_detecting_2020,mundada2019crosstalk}, oraz \glslink{qubit-leakage}{wyciek}, czyli przejście układu do stanu spoza przestrzeni obliczeniowej~\cite{Wood2015}; wskutek dwóch ostatnich błąd obwodu nie jest prostym złożeniem błędów bramek, lecz zależy od wyboru kubitów fizycznych i~od kontekstu, w~którym bramka jest wykonywana.

Skalę tych zjawisk ilustruje migawka kalibracyjna 156-kubitowego procesora ibm\_fez z~26~lutego 2025~r., dołączona do pakietu qiskit-ibm-runtime w~wersji 0.49.0~\cite{qiskit_ibm_runtime_2026}: mediany wynoszą $\Tone=145~\mu\mathrm{s}$, $\Ttwo=88~\mu\mathrm{s}$, błąd bramki jednokubitowej $2{,}3\cdot10^{-4}$, błąd bramki dwukubitowej CZ $3{,}8\cdot10^{-3}$ (bez par oznaczonych jako niesprawne) i~błąd odczytu $7{,}6\cdot10^{-3}$, więc o~jakości wyniku decyduje przede wszystkim liczba bramek dwukubitowych po transpilacji (podrozdział~\ref{subsec:tb:gates}).

Parametry szumu zmieniają się w~czasie~\cite{burnett_decoherence_2019}, podczas gdy standardowe metody charakteryzacji zakładają stałość błędów~\cite{proctor_detecting_2020}.
Według dokumentacji IBM Quantum (stan na 26~listopada 2026~r.) poszczególne parametry kalibruje się między zadaniami użytkowników, gdy wykonywane mniej więcej co godzinę zadania monitorujące wykryją odchylenie, a~raz na dobę wykonuje się pomiary wzorcowe, z~których pochodzą publikowane wartości parametrów~\cite{ibm_calibration_docs_2026}; w~latach 2021--2022, w~okresie eksperymentów z~rozdziału~\ref{chapter:pdu}, deklarowane wartości błędów bramek zmieniały się według pracy~\cite{werner_experiment_driven_2022} raz na dobę (obserwacje innych autorów: podrozdział~\ref{sec:lr:noise}).
Metoda PDU zakłada, że między kolejnymi kalibracjami błąd urządzenia jest w~przybliżeniu ustalony (założenie sprawdzane w~podrozdziale~\ref{subsec:pdu:results:stability}), a~każda charakterystyka szumu, także korekta wyznaczona eksperymentalnie, dotyczy konkretnych kubitów fizycznych i~ma ograniczony czas ważności.

\paragraph{Era NISQ i~platformy sprzętowe.}
Obecne procesory należą do okresu określanego jako \gls{nisq}~\cite{preskill_quantum_2018}: szum bramek ogranicza rozmiar obwodu, który można wykonać wiarygodnie, a~liczba kubitów nie wystarcza do pełnej korekcji błędów.
Eksperymenty w~tej pracy wykonano na procesorach IBM Quantum z~kubitami nadprzewodzącymi~\cite{krantz_quantum_2019}, w~których połączenia tylko między sąsiednimi kubitami wymuszają dodatkowe bramki dwukubitowe przy transpilacji; parametry użytych urządzeń zebrano w~dodatku~\ref{appendix:experiments}.
Kodowanie LPIQE (podrozdział~\ref{subsec:tb:qimg:lpiqe}) zaprojektowano dla układów fotonicznych, w~których fazę zadaje przesuwnik fazy, a~głównymi źródłami błędów są straty i~niepełna nierozróżnialność fotonów~\cite{slussarenko_photonic_2019,taballione_20-mode_2022}, więc wyników z~urządzeń nadprzewodzących nie można wprost przenosić na tę platformę.

\subsection{Formalizm kanałów kwantowych}\label{subsec:tb:noise:channels}

Zmiany stanu kubitu oddziałującego z~otoczeniem nie opisuje macierz unitarna, lecz \gls{quantum-channel}: liniowe odwzorowanie $\channel$ macierzy gęstości, całkowicie dodatnie i~zachowujące ślad (CPTP)~\cite{NielsenChuang2010,bib:quantum_noise}.
Każdy kanał ma postać sumy operatorowej $\channel(\rho)=\sum_{k}\kraus_k\,\rho\,\kraus_k^{\dagger}$ z~operatorami Krausa $\kraus_k$ spełniającymi $\sum_{k}\kraus_k^{\dagger}\kraus_k=\identity$.
Bramka idealna $U$ jest kanałem o~jednym operatorze Krausa, $\mathcal{U}(\rho)=U\rho U^{\dagger}$, a~bramkę zaszumioną przedstawia się jako złożenie $\errch\circ\mathcal{U}$ bramki idealnej i~\emph{kanału błędu} $\errch$.

\paragraph{Błędy koherentne i~niekoherentne.}
\Gls{coherent-error} jest unitarnym kanałem błędu, na przykład obrotem o~dodatkowy kąt $\cohangle$ wynikającym z~niedokładnej kalibracji impulsu; jest jednakowy w~każdym wykonaniu obwodu, nie zamienia stanu czystego w~mieszany i~daje się odwrócić, o~ile jest znany.
\Gls{incoherent-error}, nazywany też stochastycznym, jest kanałem, którego nie da się zapisać jednym operatorem Krausa, na przykład kanałem Pauliego, w~którym z~prawdopodobieństwem $q_j$ na stan działa macierz $\sigma_j\in\{\identity,\sigma_x,\sigma_y,\sigma_z\}$; na ogół zamienia on stany czyste w~mieszane, a~jego odwrotność nie jest kanałem.
W~błędach koherentnych sumują się amplitudy, więc $k$ powtórzeń bramki z~nadmiarowym obrotem $\cohangle$ daje prawdopodobieństwo błędnego wyniku rzędu $k^2\cohangle^2/4$, podczas gdy w~błędach niekoherentnych sumują się prawdopodobieństwa i~wzrost jest liniowy względem $k$~\cite{Wallman2015,Hashim2023}.

\paragraph{Kanały unitalne i~nieunitalne.}
Kanał jest \glslink{unital-channel}{unitalny}, jeżeli nie zmienia stanu maksymalnie mieszanego, $\channel(\identity/d)=\identity/d$, gdzie $d=2^n$; kanały unitarne i~kanały Pauliego są unitalne, a~nieunitalna jest na przykład relaksacja, która w~granicy długiego czasu każdy stan sprowadza do $\ket{0}$.
Dla jednego kubitu, w~reprezentacji~\eqref{eq:tb01:blochrho}, każdy kanał jest przekształceniem afinicznym wektora Blocha,
\begin{equation}\label{eq:tb:noise:affine}
\blochvec\;\longmapsto\;\ptmblock\,\blochvec+\ptmshift,
\end{equation}
z~rzeczywistą macierzą $\ptmblock$ o~wymiarach $3\times3$ i~wektorem przesunięcia $\ptmshift$~\cite{NielsenChuang2010}, a~unitalność oznacza $\ptmshift=\mathbf{0}$.
Dla bramki unitarnej $\ptmblock$ jest macierzą obrotu, szum niekoherentny zmniejsza jej wartości osobliwe poniżej jedności, czyli ściąga kulę Blocha, a~kanał nieunitalny przesuwa ponadto jej środek.

\subsection{Podstawowe modele szumu}\label{subsec:tb:noise:models}

Modele szumu używane w~tej pracy buduje się z~kanałów jednokubitowych~\cite{NielsenChuang2010}, których operatory Krausa oraz macierze $\ptmblock$ i~wektory $\ptmshift$ zestawia tabela~\ref{tab:tb:noise:channels} w~dodatku~\ref{app:results:tb}.
Tłumienie amplitudy o~parametrze $\qamp$ modeluje relaksację i~jest jedynym z~nich kanałem nieunitalnym, a~tłumienie fazy o~parametrze $\qphase$, które mnoży składowe $r_x$ i~$r_y$ wektora Blocha przez $c_{\varphi}=\sqrt{1-\qphase}$, modeluje czyste defazowanie; odwrócenie fazy, czyli kanał Pauliego z~macierzą $\sigma_z$ i~prawdopodobieństwem $q$, jest tym samym kanałem w~innej parametryzacji, $1-2q=c_{\varphi}$.
Złożenie obu kanałów, nazywane w~Qiskit Aer relaksacją termiczną (tu w~temperaturze zerowej), mnoży po czasie $t$ składowe $r_x$ i~$r_y$ przez $\mathrm{e}^{-t/\Ttwo}$, a~najprostszym modelem błędu koherentnego jest odstrojenie, czyli stały obrót $R_z(\cohangle)$.

Kanał depolaryzujący o~parametrze $\depolparam$ zastępuje stan $n$ kubitów mieszaniną $(1-\depolparam)\,\rho+\depolparam\,\identity/2^n$, a~dla jednego kubitu jest kanałem Pauliego, w~którym każda z~macierzy $\sigma_x$, $\sigma_y$, $\sigma_z$ działa z~prawdopodobieństwem $\qdep/3$, przy czym $\depolparam=\tfrac{4}{3}\qdep$; parametru $\depolparam$ używa Qiskit Aer~\cite{qiskit_aer_noise_docs_2026}, a~$\qdep$ PennyLane (operacja DepolarizingChannel)~\cite{bergholm_pennylane_2018}, i~nie wolno ich utożsamiać.
Po działaniu tego kanału na cały rejestr rozkład wyników jest mieszaniną rozkładu idealnego i~jednostajnego,
\begin{equation}\label{eq:tb:noise:shrink}
\probnoisy(\ket{x})=\distshrink\,\prob_{\mathrm{id}}(\ket{x})+\frac{1-\distshrink}{2^n},\qquad \distshrink=1-\depolparam,
\end{equation}
czyli odchylenia od rozkładu jednostajnego maleją w~stosunku $\distshrink$, nazywanym dalej \glslink{distribution-shrinkage}{współczynnikiem skurczenia rozkładu}.
Błąd odczytu opisuje się klasycznie, jako losową zamianę etykiet wyników~\cite{maciejewski_mitigation_2020}:
\begin{equation}\label{eq:tb:noise:readout}
\probnoisy(\ket{y})=\sum_{x\in\{0,1\}^n}\confmat_{yx}\,\prob(\ket{x}),\qquad \confmat_{yx}=q_{y|x},
\end{equation}
gdzie $\prob$ jest rozkładem przed odczytem, a~$q_{y|x}$ prawdopodobieństwem odczytania wyniku $y$ dla rejestru w~stanie $\ket{x}$; wyznaczanie i~odwracanie macierzy $\confmat$ opisano w~podrozdziale~\ref{subsec:tb:qec:mitigation}.

\paragraph{Model warstwowy.}
W~rozdziałach~\ref{chapter:learned_encoding} i~\ref{chapter:classification} używany jest fenomenologiczny model, w~którym po każdym bloku obwodu, czyli powtarzanym fragmencie ansatzu (w~kodzie: warstwie, stąd nazwa), na każdy kubit działają odstrojenie $R_z(\cohangle)$, tłumienie fazy o~parametrze $\qphase$ i~kanał depolaryzujący o~parametrze $\qdep$.
Symulacje z~tym modelem nie zawierają błędu odczytu~\eqref{eq:tb:noise:readout} (zawiera go tylko model budowany z~danych kalibracyjnych, opisany niżej), a~w~części konfiguracji wyłączone jest tłumienie fazy albo odstrojenie.
Te trzy odwzorowania są przemienne, więc ich kolejność w~obrębie bloku nie wpływa na wynik.
W~eksperymencie interferencyjnym z~przykładu~\ref{ex:tb01:interference} jedna dawka takiego szumu przed drugą bramką Hadamarda daje
\begin{equation}\label{eq:tb:noise:fringe}
\probnoisy(\ket{0})=\tfrac{1}{2}\bigl[1+v\cos(\varphi+\cohangle)\bigr],\qquad v=c_{\varphi}\bigl(1-\tfrac{4}{3}\qdep\bigr),
\end{equation}
czyli składowa koherentna przesuwa argument funkcji, a~składowe niekoherentne zmniejszają widoczność interferencji $v$ (podrozdział~\ref{subsec:tb:density}).

\paragraph{Szum a~składowe błędu wyniku.}
W~rozkładzie~\eqref{eq:tb01:errordecomp} błąd sprzętowy $\probnoisy-\prob_{\mathrm{id}}$ jest w~całości wyznaczony przez kanał obwodu i~macierz odczytu, a~wchodzą do niego zarówno błędy koherentne, jak i~niekoherentne: w~równaniu~\eqref{eq:tb:noise:fringe} przesunięcie $\cohangle$ i~widoczność $v$ łącznie określają deterministyczną funkcję fazy, powtarzalną, dopóki parametry urządzenia pozostają ustalone.
Podział na błędy koherentne i~niekoherentne nie pokrywa się zatem z~podziałem na składową systematyczną i~statystyczną, która wynika wyłącznie ze skończonej liczby pomiarów.
Rozróżnienie wyznacza granice klasycznej korekty wyniku (rozdziały~\ref{chapter:pdu}, \ref{chapter:learned_encoding} i~\ref{chapter:classification}): korekta addytywna, czyli odjęcie oszacowania składowej systematycznej, nie zmienia wariancji składowej statystycznej poza niepewnością samego oszacowania, natomiast korekta odwracająca skurczenie~\eqref{eq:tb:noise:shrink} wymaga dzielenia przez $\distshrink<1$, a~zatem zwiększa odchylenie standardowe składowej statystycznej $1/\distshrink$ razy.

\paragraph{Modele szumu w~symulatorach.}
W~modelu budowanym w~Qiskit Aer z~danych kalibracyjnych urządzenia po każdej bramce działa kanał depolaryzujący, a~po nim relaksacja termiczna o~parametrach $\Tone$, $\Ttwo$ i~czasie trwania bramki, przy czym parametr depolaryzacji dobiera się tak, aby średni błąd złożenia był równy błędowi bramki z~kalibracji; odczyt opisują niezależne jednokubitowe macierze $\confmat$~\cite{qiskit_aer_noise_docs_2026}.
Model ten nie zawiera błędów koherentnych, przesłuchu, wycieku, korelacji błędów odczytu ani dryfu, a~parametry dopasowane do wyników urządzenia różnią się od kalibracyjnych~\cite{georgopoulos_modeling_2021}, dlatego symulacja z~szumem jest przybliżeniem, a~nie zamiennikiem pomiaru na sprzęcie.

\subsection{Charakteryzacja szumu}\label{subsec:tb:noise:characterisation}

Tomografia procesu kwantowego (QPT) odtwarza kanał z~tomografii stanów wyjściowych, czyli z~pomiarów w~kilku bazach, dla zbioru stanów wejściowych rozpinającego przestrzeń macierzy gęstości~\cite{bib:quantum_tomography,ChuangNielsen1997,Poyatos1997}.
Dla jednego kubitu wystarczają cztery stany wejściowe i~trzy bazy pomiarowe, czyli dwanaście obwodów, lecz liczba parametrów kanału rośnie jak $16^n$~\cite{NielsenChuang2010,MohseniQPTReview}, a~badanej operacji przypisuje się także błędy SPAM~\cite{Merkel2013}.
Wynik QPT przedstawia się jako \gls{ptm} (PTM)~\cite{greenbaum2015introduction}, która dla jednego kubitu działa na wektor $(1,r_x,r_y,r_z)^{\top}$ i~ma postać blokową:
\begin{equation}\label{eq:tb:noise:ptm-blocks}
R_{ab}=\frac{1}{2}\Tr\bigl[\sigma_a\,\channel(\sigma_b)\bigr],\quad a,b\in\{0,x,y,z\},\quad \sigma_0=\identity,
\qquad
\ptm=\begin{pmatrix} R_{00} & \ptmrow^{\top}\\ \ptmshift & \ptmblock \end{pmatrix},
\end{equation}
gdzie $\ptmblock$ i~$\ptmshift$ pochodzą z~równania~\eqref{eq:tb:noise:affine}.
Macierz $\ptmblock$ nazywa się \emph{blokiem unitalnym} (dla bramki $\mathrm{X}$ jest równa $\operatorname{diag}(1,-1,-1)$), a~wektor $\ptmshift$ częścią nieunitalną kanału~\cite{greenbaum2015introduction,Wallman2015}.
Dla kanału zachowującego ślad $R_{00}=1$ i~$\ptmrow=\mathbf{0}$; niezerowe elementy $\ptmrow$ są możliwe tylko w~rekonstrukcji bez narzuconego zachowania śladu~\cite{Wallman2015}.

Błędy bramek w~danych kalibracyjnych wyznacza się testem z~losowymi sekwencjami bramek (ang. \emph{randomized benchmarking}, RB)~\cite{Magesan2011,magesan2012characterizing,ibm_calibration_docs_2026}, który przy błędzie niezależnym od bramki uśrednia kanał błędu do kanału depolaryzującego o~tej samej średniej wierności; uzasadnia to użycie tego kanału w~modelach symulatorów, ale sprawia też, że wynik RB nie rozróżnia błędów koherentnych i~niekoherentnych~\cite{Wallman2015,ProctorWhat2017}.
